\documentclass[10pt,twoside]{article}
\usepackage[utf8]{inputenc}
\usepackage{amsmath}
\usepackage{amsfonts}
\usepackage{amssymb}
\usepackage[spanish,es-noshorthands]{babel}
\usepackage[T1]{fontenc}
\usepackage{lmodern}
\usepackage{graphicx,hyperref}
\usepackage{tikz,pgf}
\usepackage{multicol}
\usepackage{subfig}
\usepackage[papersize={6.5in,8.5in},width=5.5in,height=7in]{geometry}
\usepackage{fancyhdr}
\pagestyle{fancy}
\fancyhead[LE]{\includegraphics[height=12pt]{Images/logo-colegio.png} \'Algebra $8^{\circ}$}
\fancyhead[RE]{}
\fancyhead[RO]{\textit{Germ\'an Avenda\~no Ram\'irez, Lic. U.D., M.Sc. U.N.}}
\fancyhead[LO]{}

\author{Germ\'an Avenda\~no Ram\'irez, Lic. U.D., M.Sc. U.N.}
\title{\begin{minipage}{.2\textwidth}
\includegraphics[height=1.75cm]{Images/logo-colegio.png}\end{minipage}
\begin{minipage}{.55\textwidth}
\begin{center}
Taller 02, Introducción al álgebra\\
Álgebra $8^{\circ}$
\end{center}
\end{minipage}\hfill
\begin{minipage}{.2\textwidth}
\includegraphics[height=1.75cm]{Images/logo-sed.png} 
\end{minipage}}
\date{}
\thispagestyle{plain}
\begin{document}
\maketitle
Nombre: \hrulefill Curso: \underline{\hspace*{44pt}} Fecha: \underline{\hspace*{2.5cm}}
 \section*{Nivel I}
 \begin{enumerate}
 \item Expresa, indicando las operaciones que debes hacer y calculando el resultado: 
\begin{center}
\begin{tabular}{|l|c|c|}
\hline 
\hspace*{.5cm}Expresión & Resultado & Operación indicada \\ 
\hline 
El doble de 5 es \ldots & 10 & $2\cdot 5$ \\ 
\hline 
La mitad de 8 es \ldots &  &  \\ 
\hline 
El triple de 9 es \ldots &  &  \\ 
\hline 
El cuadrado de 7 es \ldots &  & \\ 
\hline 
La raíz cuadrada de 25 es \ldots &  &  \\ 
\hline 
La suma de 8 y 5 es \ldots &  &  \\ 
\hline 
La diferencia entre 10 y 7 es \ldots &  &  \\ 
\hline 
El producto de 4 y 10 es \ldots &  &  \\ 
\hline 
El cociente entre 24 y 8 es \ldots &  &  \\ 
\hline 
\end{tabular} 
\end{center}
\end{enumerate}



\end{document}
